\scanpage{24}
\(2^\circ)\) Un hom \(\pi\xrightarrow{\varphi'}\widetilde G'_0\) qui centralise \(N\) et remonte \(\pi\to\unsure{\mathfrak{S}'|\mathcal{G}'|\mathfrak{G}'}\) [donc, si \(\pi\subset\widetilde H_0\) d'où \(\deleted{\pi\to}\pi\to\mathfrak{S}'\lhook\joinrel\to\mathfrak{S}\) trivial, un hom
\[
 \pi\xrightarrow{\varphi'}\mathfrak{z}(N)\qquad ]
\]

Ceci posé, on vérifie immédiatement que l'extension \(\widetilde G'_0\) est l'image inverse de \(\widetilde G_0\) par \(\mathfrak{S}'\to\mathfrak{S}\), \deleted{et que} (donc \(\widetilde G'_0\subset\widetilde G_0\)), et que \(\varphi'\) est induit par \(\varphi\) [donc si \(\pi\subset\widetilde H_0\), \(\varphi=\varphi'\) en tant que hom \(\pi\to\mathfrak{z}(N)\)]

Si les conditions \deleted{\unsure{\(\alpha)\)}} \(\beta)\ \gamma)\) sont \unsure{vérifiées|vérifiés} pour les données \deleted{\unsure{\(a)\)}}\inserted{\(a)\)} \(b)\ \unsure{e|c})\), il en est de \(=\) pour les données \deleted{\unsure{\(a)\ b)\ e)\)}} \(a')\,b')\,e')\) (\(\pi,N,\mathfrak{S}'\) discrets, \(\pi\subset\widetilde H'_0\)) mais la condition \(\alpha\) (\(\widetilde H_0^{\deleted{\unsure{0|\prime|\ast}}}\) connexe et simplement connexe n'implique pas néc. \(=\) que \(\widetilde H'_0{}^{\deleted{\unsure{0}}}\) soit connexe \ldots)

\deleted{Mais comme \(\mathfrak{S}'\) est discret, \(\deleted{\unsure{\widetilde H'_0}}\), il en est de \(=\) de \(G'_0\), donc \(G'_0\) apparaît comme un quotient de \(\pi_0\widetilde G'\)}
% archive sheet number: 23

